\section{Asymptotic closure}
\label{sec:asym}

If a small parameter makes the discarded variables relax quickly to a state
determined by the retained ones, the closure can be \emph{derived}, with an
error estimate. This is the only kind of closure that comes with one.

\subsection{The small parameter and two expansions}

For a gas the parameter is the Knudsen number
\begin{equation}
  \Kn = \frac{\lambda}{L},
\end{equation}
the mean free path over the gradient length. In units of the macroscopic
length and time the Boltzmann equation reads
$\partial_tf+\bv\cdot\nabla f=\Kn^{-1}C[f]$. As $\Kn\to0$ the collision
term forces $C[f]=0$, whose solutions are the local Maxwellians
\[
  f_M(n,\bu,T) = n\Bigl(\frac{m}{2\pi T}\Bigr)^{3/2}
  \exp\Bigl(-\frac{m|\bv-\bu|^2}{2T}\Bigr).
\]
Two expansions organize the corrections.
\begin{itemize}
  \item \textbf{Hilbert (1912)} expands both $f$ and the fluid fields in
    powers of $\Kn$. Every order solves linearized Euler equations forced by
    the previous orders. The result is a series solution, not a new set of
    equations.
  \item \textbf{Chapman--Enskog (1916--17)} expands $f$ but not the fluid
    fields:
    $f=f_M(n,\bu,T)\,(1+\Kn\,\phi_1+\Kn^2\phi_2+\dots)$, with each
    $\phi_j$ a function of the fields and their gradients and carrying no
    density, momentum or energy. The time derivative is expanded instead.
    The result is a sequence of closed equations for $n$, $\bu$, $T$: Euler,
    Navier--Stokes--Fourier, Burnett.
\end{itemize}

\subsection{First order: Navier--Stokes--Fourier}
\label{sec:ce}

Let $\Lambda$ be the linearized collision operator about $f_M$, written so
that $C[f_M(1+\phi)] = -f_M\,\Lambda\phi+O(\phi^2)$. It is self-adjoint and
non-negative in the Maxwellian-weighted inner product
$\avg{\phi,\psi}=n^{-1}\int\phi\psi f_M\,\dd^3v$, and its kernel is spanned
by the collision invariants $1$, $\bv$, $|\bv|^2$.

At order $\Kn^0$ the kinetic equation gives
\begin{equation}
  \Lambda\phi_1 = -\bigl(\partial_t+\bv\cdot\nabla\bigr)\ln f_M ,
  \label{eq:ce1}
\end{equation}
where the time derivatives on the right are replaced by the Euler
equations. Writing $\bw=\bv-\bu$, a direct computation gives
\begin{equation}
  \bigl(\partial_t+\bv\cdot\nabla\bigr)\ln f_M
  = \Bigl(\frac{m|\bw|^2}{2T}-\frac52\Bigr)\bw\cdot\nabla\ln T
  + \frac{m}{T}\Bigl(\bw\bw-\tfrac13|\bw|^2\mathbf I\Bigr):\nabla\bu .
  \label{eq:ce-source}
\end{equation}
Three points deserve attention.
\begin{enumerate}
  \item \emph{Solvability.} Equation \eqref{eq:ce1} can be solved only if
    its right-hand side is orthogonal to the kernel of $\Lambda$. The Euler
    equations are exactly that condition. This is the Fredholm alternative,
    and it is where the fluid equations come from at every order.
  \item \emph{The source has two pieces,} a vector multiplying $\nabla T$
    and a traceless tensor multiplying $\nabla\bu$. By isotropy of
    $\Lambda$ the first produces only a heat flux and the second only a
    stress. No density-gradient term survives.
  \item \emph{The transport coefficients are quadratic forms.} With
    dimensionless $\bc=\bw/\vt$, $\vt=\sqrt{T/m}$, and
    \[
      \mathbf A=\bigl(\tfrac12|\bc|^2-\tfrac52\bigr)\bc,\qquad
      \mathbf B=\bc\bc-\tfrac13|\bc|^2\mathbf I,
    \]
    the fluxes are
    \begin{equation}
      \boldsymbol{\pi} = -\mu\left(\nabla\bu + \nabla\bu^{\mathsf T}
         - \tfrac{2}{3}(\nabla\cdot\bu)\,\mathbf{I}\right),
      \qquad
      \bq = -\kappa\,\nabla T ,
      \label{eq:NSF}
    \end{equation}
    \begin{equation}
      \kappa = \frac{nT}{m}\cdot\frac13\avg{\mathbf A,\Lambda^{-1}\mathbf A},
      \qquad
      \mu = nT\cdot\frac1{10}\avg{\mathbf B,\Lambda^{-1}\mathbf B} .
      \label{eq:transport-forms}
    \end{equation}
\end{enumerate}
Because $\Lambda$ is non-negative, each coefficient is the maximum of a
concave quadratic functional,
\begin{equation}
  \avg{\mathbf A,\Lambda^{-1}\mathbf A}
  = \max_{\boldsymbol\varphi}\bigl\{2\avg{\mathbf A,\boldsymbol\varphi}
      -\avg{\boldsymbol\varphi,\Lambda\boldsymbol\varphi}\bigr\},
  \label{eq:variational}
\end{equation}
so any trial function gives a lower bound. This is how transport
coefficients have been computed since Chapman and Enskog (polynomial trial
functions, the Sonine expansion), and it is the starting point of the
two-sided bounds proved in this repository (Section~\ref{sec:repo}).

\begin{example}[BGK]
\label{ex:bgk}
The BGK model replaces collisions by relaxation to the local Maxwellian at
rate $\nu$, $C[f]=\nu(f_M-f)$ \cite{BGK54}. Then $\Lambda=\nu$ on the
complement of the collision invariants, and \eqref{eq:transport-forms}
needs only Gaussian moments:
$\tfrac13\avg{\mathbf A,\mathbf A}=\tfrac52$ and
$\tfrac1{10}\avg{\mathbf B,\mathbf B}=1$. Hence
\begin{equation}
  \kappa=\frac52\,\frac{nT}{m\nu},\qquad \mu=\frac{nT}{\nu},\qquad
  \Pr=\frac{c_p\,\mu}{\kappa}=1 \quad(c_p=\tfrac52/m).
\end{equation}
A monatomic gas has $\Pr\approx\tfrac23$. With one relaxation rate BGK
cannot match the viscosity and the conductivity together, which is a first
instance of a theme of Section~\ref{sec:between}: a model collision
operator fitted to one transport coefficient does not reproduce the others.

In one velocity dimension, with the heat flux defined as the third central
moment (the convention of Section~\ref{sec:collisionless}), the same
calculation gives $\kF=3\,nT/(m\nu)$.
\end{example}

For hard spheres of diameter $d$ the first Sonine approximation gives
$\mu=\tfrac{5}{16d^2}\sqrt{mT/\pi}$ and $\Pr=\tfrac23$; higher
approximations change $\mu$ and $\kappa$ by about two percent
\cite{ChapmanCowling}. In both cases
$\mu\sim mn\,\vt\lambda$ and $\kappa\sim n\,\vt\lambda$: a mean free path
times a thermal speed.

\subsection{Higher order, and why it fails}
\label{sec:burnett}

The second order of the Chapman--Enskog expansion gives the Burnett
equations, with terms of second order in gradients. They are not an
improvement in practice. Bobylev showed that they are linearly unstable at
short wavelengths \cite{Bobylev82}.

The mechanism can be seen in one line. For a linear temperature
perturbation $\propto e^{ikx}$ write the heat flux as
$\hat q=-ik\,\chi(k)\,\hat T$, so that the temperature equation is
$\partial_t\hat T\propto-k^2\chi(k)\hat T$. Fourier's law is
$\chi=\kF$. The expansion continued in powers of the gradient gives
\begin{equation}
  \chi(k)=\kF\bigl(1-K\,k^2+O(k^4)\bigr),\qquad K>0,
  \label{eq:deficit}
\end{equation}
where the sign reflects that the true flux falls below Fourier's as
gradients steepen. Truncating after the $k^2$ term makes $\chi$ negative for
$k^2>1/K$, where the temperature mode grows at a rate proportional to
$k^4$: an ill-posed equation. The rational form
\begin{equation}
  \chi(k)=\frac{\kF}{1+K\,k^2}
  \label{eq:rational}
\end{equation}
has the same expansion to this order and is positive and bounded at every
$k$. This is the same distinction as in Example~\ref{ex:fastslow} between
the polynomial and the rational first-order closures. For conserving
collision operators in one velocity dimension, \eqref{eq:rational} is a
proved lower bound on the exact response, and $\kF$ is an upper bound
(Section~\ref{sec:repo}).

Repairs of the Burnett equations take three routes: regularized moment
systems such as R13 \cite{StruchtrupTorrilhon03,Struchtrup05}; exact
summation of the expansion in linear problems, and invariant-manifold
constructions \cite{GorbanKarlin14,KogelbauerKarlin24}; and abandoning the
expansion for a moment closure (Section~\ref{sec:moment}).

\subsection{What is proved}

The formal expansions are backed by theorems in several settings. Caflisch
justified the compressible Euler limit for smooth solutions
\cite{Caflisch80}. Bardos, Golse and Levermore set out the program for
incompressible limits \cite{BGL91}, and Golse and Saint-Raymond carried it
through for the incompressible Navier--Stokes limit in the setting of
renormalized solutions \cite{GSR04}. Deng, Hani and Ma connect hard-sphere
dynamics through the Boltzmann equation to the fluid equations
\cite{DHM25}. These are the theorem-faithful settings that Track~A of the
benchmark uses.

\subsection{Why neutral gases are the easy case}
\label{sec:easy}

Collisions between neutral molecules are short-range and frequent. Three
facts follow, and each fails for plasmas (Section~\ref{sec:plasma}).
\begin{enumerate}
  \item \textbf{Fast local relaxation.} Collisions drive $f$ to a local
    Maxwellian in a few collision times, so the kinetic state is determined
    by $n$, $\bu$, $T$ and their gradients.
  \item \textbf{One small parameter.} All molecules have comparable mean
    free paths, so a single $\Kn$ controls the expansion.
  \item \textbf{Locality.} Fluxes depend on gradients at the same point and
    time.
\end{enumerate}
